% !TeX root = ../main.tex
\begin{englishabstract}

With the rapid development of wearable devices and artificial intelligence technologies, modeling human motion states using non-visual sensors has become an increasingly active research topic. Among various sensing modalities, inertial measurement units (IMUs) can capture high-frequency kinematic information of human motion, while channel state information (CSI) extracted from WiFi signals can reflect disturbances in the wireless propagation environment caused by human activities. Compared with vision-based approaches, these sensing modalities are more privacy-preserving, flexible to deploy, and robust to environmental conditions. Nevertheless, several challenges remain in practical applications. First, heterogeneous sensing modalities exhibit significant differences in signal forms and feature distributions, making effective multimodal alignment and fusion difficult. Second, variations in movement patterns and physical characteristics across individuals often lead to performance degradation in cross-subject scenarios. In addition, existing inertial-sensor-based human pose reconstruction methods mainly focus on estimating poses at the current time step, while lacking the capability to model future poses.

To address these challenges, this thesis investigates two tasks: human activity recognition and human pose reconstruction. For the activity recognition task, a multimodal alignment and fusion based method named WIHAR is proposed to tackle the difficulties of heterogeneous feature alignment between IMU and WiFi CSI data and the limited generalization ability in cross-subject scenarios. Specifically, a cross-modal contrastive learning strategy \cite{radford2021learning} is employed to map IMU motion dynamics and WiFi CSI radio-frequency sensing features into a unified semantic space. A cross-attention mechanism is further introduced to enable feature interaction between the two modalities. Moreover, a prototype-based few-shot adaptation strategy is adopted to alleviate the feature distribution discrepancy between source and target subjects by aligning target subject features with class prototypes learned from the source domain. Experimental results on the XRFV2 dataset \cite{lan2025xrf} show that WIHAR achieves an activity recognition accuracy of 96.06\%. In cross-subject experiments, the initial recognition accuracy of the model is 72.57\%, which can be improved to 90.99\% after fine-tuning with a small amount of data from the target subject.

For the human pose reconstruction task, an autoregressive modeling mechanism based method named AIPose is proposed to address the lack of temporal prediction capability in existing approaches. The proposed method takes IMU data as the only input and constructs a closed-loop architecture consisting of an IMU feature extraction module, a pose reconstruction module, and an IMU prediction module. By predicting future IMU sequences from the current IMU sequence, the model further performs pose prediction for future time steps. In this way, the framework enables both current pose reconstruction and future pose prediction simultaneously. Experimental results on the XRFV2 dataset show that AIPose achieves a mean per-joint position error (MPJPE) of 50.56 for current pose reconstruction and 56.29 for future pose prediction.

\englishkeywordstype{Activity recognition; Pose reconstruction; Multimodal fusion; Cross-modal alignment; Autoregressive modeling}{Application Research}

\end{englishabstract}